\section{The relation quotient}
\label{sec:relation}

Let $V_n=\QQ[X_n]$ and define
\[
  W_n=\im(I+S)+\im(I+R+R^2),\qquad M_n=V_n/W_n .
\]
Let $o_S(n)$ and $o_R(n)$ denote the numbers of $S$- and $R$-orbits on $X_n$.
Construct a bipartite multigraph $D_n$: left vertices are $S$-orbits, right
vertices are $R$-orbits, and each state $x\in X_n$ is an edge connecting the
two orbit vertices containing $x$.

\begin{theorem}[Exact Manin-rank ledger]
\label{thm:rank}
For every $n\ge2$,
\[
  \dim_\QQ M_n=|X_n|-o_S(n)-o_R(n)+1=\beta_1(D_n).
\]
\end{theorem}

\begin{proof}
Over $\QQ$, the image of $I+S$ is spanned by the indicator vectors of
$S$-orbits, and the image of $I+R+R^2$ is spanned by the indicator vectors of
$R$-orbits.  If a vector lies in both spans, its coefficient on a state depends
only on the $S$-orbit and only on the $R$-orbit.  The graph $D_n$ is connected
because the $S,R$ action on $P^1(\ZZ/n\ZZ)$ is transitive.  Hence the common
intersection is the all-ones line.  Therefore
\[
  \dim W_n=o_S(n)+o_R(n)-1,
\]
and the displayed formula follows.  The same expression is exactly the
cycle-rank formula for the connected graph $D_n$.
\end{proof}

\begin{remark}
Equivalently, $M_n=H^1(D_n;\QQ)$ after orienting each state edge from its
$S$-orbit vertex to its $R$-orbit vertex.  This is why the quotient naturally
lands in modular-symbol homology.  The interpretation is useful, but it is not
a new determinant object.
\end{remark}

\begin{corollary}[Cuspidal quotient is not a selector]
\label{cor:cuspidal}
Removing Eisenstein cusp-boundary directions cannot rescue prime selectivity
inside this family.
\end{corollary}

\begin{proof}[Audit-supported proof]
Classically, the stronger cuspidal dimension is $2g_0(n)$.  The exact audit
through $n=192$ finds nonzero cuspidal homology on 139 of 148 composite blocks
and zero cuspidal homology on five prime blocks.  It therefore fails both
directions of a field selector.
\end{proof}
